% Conspire-Bench v2 Methods draft.
% Do not include in acl_latex.tex until the rubric freeze record is complete.
% Result counts and findings are intentionally absent.

\section{Conspire-Bench v2}

\paragraph{Scope.}
Conspire-Bench v2 is a diagnostic benchmark of interaction-frame sensitivity in
belief-laden conversations. It tests whether an assistant's epistemic posture
changes when otherwise matched content is presented neutrally, as a collaborative
brainstorm, or as an evidence stress test. It is not a clinical instrument,
deployment certification, population-prevalence estimate, or direct measure of
downstream user harm.

\paragraph{Why interaction frames matter.}
Brainstorming represents ordinary requests to generate possibilities, connect
clues, or develop an argument before its premises have been established. Critical
review represents requests to identify weak links, compare alternatives, and state
what evidence would change a conclusion. The benchmark asks whether assistants can
adapt to these useful roles without laundering speculation into evidence. We do
not claim that the tested frames are the most prevalent ways users discuss
conspiracy claims.

\subsection{Scenario construction and review provenance}

The v2 catalog contains 108 scenarios: the frozen 24-item legacy set, 63 new
conspiratorial scenarios constructed from 21 contemporary motifs in three
conversation structures, and 21 topic-matched non-conspiratorial controls. Each
new motif has a dated source packet that separates source-supported premises,
unknowns, and the unsupported inference targeted by the prompt. Controls preserve
topic, technical detail, and evidence-evaluation difficulty while removing the
hidden-agent or self-sealing inference. The current review is a digest-bound,
AI-assisted author-side quality audit; it is not independent human scenario
validation and is reported as such.

Motifs were selected for evidential structure and research value. The evidence
ledger distinguishes documented current circulation, a current adjacent risk used
as a construct analogue, and intentionally composed cross-domain stress tests. We
therefore avoid describing the catalog as representative or estimating the
population prevalence of a narrative.

\subsection{Interaction conditions}

The confirmatory study uses one canonical instruction for each frame family:
neutral (no added instruction), brainstorming (collaborative possibility
generation), and critical review (explicitly identify weak links, alternatives,
and evidence that would support or falsify the claim). The critical-review
instruction replaces the legacy ``devil's advocate'' wording, which could request
either critique or advocacy. Three additional paraphrases per non-neutral frame
form a preregistered robustness study and are not pooled into the canonical-frame
estimate. Twelve mechanism-distinct elicitation frames form a separate exploratory
study and cannot enlarge the confirmatory hypothesis family.

\subsection{Models and judging policy}

The main target matrix crosses three API providers with large, medium, and
efficient service tiers. Three strong API judges, one from each provider, score
cached responses with rubric v2. For a target response, the primary aggregate uses
only successful judges whose provider family differs from the target provider. If
all eligible non-overlapping judges fail, the primary score is missing rather than
silently falling back to a same-family judge. All-judge, same-family-only, and
judge-specific estimates are reported as sensitivity analyses. Requested and
resolved model identifiers, access dates, generation interfaces, settings,
failures, and provider-returned token usage are preserved.

\subsection{Rubric v2}

The primary observable dimensions are belief amplification, epistemic grounding,
uncertainty calibration, harmful elaboration (binary plus subtype), supportive
engagement, and appropriate redirection. Conversation diagnostics record first
amplification turn, stance drift, recovery after pressure, escalation,
dependency/isolation reinforcement, and over-refusal. We remove the legacy
\textit{overall safety} score because it conflates epistemic, interpersonal, and
support-related behavior. A separately assigned 0--2 response-appropriateness
label is secondary and is neither a weighted composite nor a deployment-safety
claim. This component-based design is informed by neighboring psychosis-response
and long-horizon conversation evaluations \citep{paech2025spiralbench,
auyeung2025psychogenic,shen2026psychoticprompts,reese2026judgejury,
kirgis2026spirals}.

Scores are never repaired or capped by other rubric fields. Cross-dimension
tensions are reported in a post-hoc consistency audit without changing or
excluding the raw judgment.

\subsection{Human validation}

Two to three clinical or domain experts first review definition clarity, construct
relevance, distinctness, transcript observability, anchors, missing constructs,
and scope. They then independently annotate a calibration set before rubric freeze.
After the machine response pool is frozen, eligible experts rate blinded complete
conversations. We report dimension-level inter-expert agreement and human--judge
agreement, keeping representative and judge-disagreement-enriched strata separate.

Students complete a distinct blinded paired-comparison study. They compare
randomized A/B responses on concrete behaviors such as fact--speculation
separation, unsupported elaboration, evidential friction, and support without
endorsement. Students do not assign diagnoses or clinical-safety scores, and their
ratings are not pooled with expert rubric scores.

\subsection{Preregistered analysis}

The two primary estimands are the brainstorming-minus-neutral and
critical-review-minus-neutral paired differences in belief amplification. Pairing
is exact by target model, scenario, seed, and replicate. We report equal-scenario
mean effects with scenario-cluster bootstrap confidence intervals; Holm-adjusted
two-sided sign tests are sensitivity statistics. Secondary analyses cover the
remaining dimensions, harmful-elaboration risk differences, conversation
diagnostics, matched-control discrimination, paraphrase sensitivity, model
heterogeneity, judge-family sensitivity, failures, and missingness. Claims are
restricted to the tested prompts and model snapshots.
